\documentclass[../../../../main.tex]{subfiles}
\begin{document}

{}
　$i$を虚数単位とする．実数$\theta$に対して複素数$e(\theta)$を
\[ e(\theta)=\cos\theta+i\sin\theta \]
で定める．正の整数$n$に対し，複素数平面上の
点$\displaystyle e \left( \frac{n-1}{3}\pi \right)$と
点$\displaystyle e \left( \frac{n}{3}\pi \right)$を通る直線を$l_n$とする．
複素数$z$に対し，直線$l_n$に関して点$z$と対称な位置にある点を表す複素数を$R_n(z)$とする．
\begin{description}
\item[(1)]複素数$z$と共役な複素数を$\overline{z}$とする．このとき，
\[ R_n(z)=e(\theta_1)\overline{z}+\sqrt3 \, e(\theta_2) \]
をみたす実数$\theta_1, \, \theta_2$をそれぞれ$n$を用いて表せ．
\item[(2)]実数$a, \, b$に対して，
\[ z_1=a+bi, \hspace{1zw} z_{n+1}=R_n(z_n) \hspace{1zw} (n=1, \, 2, \, 3, \, \cdots) \]
で定められる複素数の列$\{ z_n \}$の一般項を求めよ．
\end{description}

\end{document}
